% !TEX root = master.tex
% !TEX program = lualatex

\chapter{Pointers - 1}%
\label{cha:Pointers - 1}

\section{C04a — Introduction aux pointeurs : à quoi servent-ils?}

\subsection{Démystifier les pointeurs}

Les pointeurs sont souvent perçus comme un concept difficile en C.
En réalité, leur idée fondamentale est simple.

Un pointeur est simplement :

\begin{itemize}
    \item une \textbf{référence vers une donnée en mémoire}
    \item un moyen d’accéder à un objet sans le copier
\end{itemize}

On peut comparer un pointeur à un \textbf{lien} ou un \textbf{signet} :

\begin{itemize}
    \item quand on partage un site web, on ne copie pas tout le site
    \item on partage seulement un lien vers ce site
\end{itemize}

De la même manière, un pointeur est un lien vers une donnée.
Il permet d’accéder à cette donnée depuis différents endroits du programme
sans la dupliquer.

\subsection{Les trois usages fondamentaux des pointeurs}

Les pointeurs servent principalement à trois choses.

\subsubsection{1. Partage d’un objet (référence)}

On peut vouloir utiliser le même objet à plusieurs endroits du programme.

\begin{itemize}
    \item éviter la duplication de données
    \item partager une structure entre fonctions
    \item permettre à plusieurs fonctions de modifier le même objet
\end{itemize}

Dans ce cas, le pointeur agit comme une \textbf{référence vers l’objet}.
On manipule toujours le même objet en mémoire.

\subsubsection{2. Généricité}

Parfois, on ne sait pas à l’avance quel objet on veut manipuler.

Un pointeur permet alors :

\begin{itemize}
    \item de décrire un objet \textbf{de façon générique}
    \item de choisir dynamiquement l’objet utilisé
    \item d’écrire du code réutilisable
\end{itemize}

\paragraph{Exemple : pointeurs vers fonctions}

Supposons que l’on possède plusieurs fonctions mathématiques :

\[
f_1(x),\; f_2(x),\; f_3(x),\; \dots
\]

On veut écrire une fonction qui calcule une intégrale,
mais la fonction à intégrer est choisie par l’utilisateur.

Sans pointeurs :

\begin{itemize}
    \item il faudrait recompiler le programme pour changer la fonction
\end{itemize}

Avec un pointeur vers fonction :

\begin{itemize}
    \item on choisit dynamiquement la fonction
    \item on passe un pointeur vers la fonction choisie
\end{itemize}

Le pointeur devient alors une description générique
de la fonction à manipuler.

\subsubsection{3. Allocation dynamique (contrôle de la durée de vie)}

Les pointeurs permettent aussi de gérer la durée de vie des objets.

On distingue deux notions importantes :

\begin{itemize}
    \item \textbf{Portée (scope)} :
    ensemble des lignes où une variable est visible dans le code
    \item \textbf{Durée de vie} :
    intervalle de temps où l’objet existe réellement en mémoire
\end{itemize}

Pour une variable locale classique :

\begin{itemize}
    \item elle est créée au début du bloc
    \item elle est détruite à la fin du bloc
\end{itemize}

Sa durée de vie correspond donc à la durée d’exécution de sa portée.

Cependant, on peut vouloir :

\begin{itemize}
    \item créer un objet dans une fonction
    \item continuer à l’utiliser après le retour de la fonction
    \item contrôler manuellement sa création et sa destruction
\end{itemize}

Dans ce cas, on utilise l’\textbf{allocation dynamique}
(qui sera étudiée en détail plus tard).

Les pointeurs permettent alors d’accéder à cet objet
au-delà de sa portée syntaxique.

\section{C04b — Représentation concrète des pointeurs}

\subsection{Un pointeur est une variable}

En C, un pointeur est une variable comme les autres.

\begin{itemize}
    \item il occupe un emplacement en mémoire
    \item il possède une valeur stockée
\end{itemize}

La seule différence est que cette valeur n’est pas un nombre ordinaire,
mais une \textbf{adresse mémoire}.

Autrement dit :

\begin{center}
\textbf{un pointeur stocke l’adresse d’une autre variable}
\end{center}

On peut voir un pointeur comme une
\textbf{variable qui contient l’adresse d’une variable}.

Cette adresse permet d’accéder à la zone mémoire correspondante
et donc à la valeur de l’objet pointé.

\subsection{Analogie : carnet d’adresses}

Une bonne façon de comprendre les pointeurs
est de les voir comme des pages dans un carnet d’adresses.

\begin{itemize}
    \item une \textbf{maison} représente une donnée en mémoire
    \item une \textbf{adresse} correspond à l’emplacement mémoire
    \item une \textbf{page du carnet} correspond à un pointeur
\end{itemize}

Le pointeur ne contient pas la maison,
il contient seulement son adresse.

\paragraph{Conséquence importante}

Déclarer un pointeur ne crée pas l’objet pointé.

\begin{itemize}
    \item on ajoute simplement une nouvelle page dans le carnet
    \item cette page peut être vide
    \item elle peut contenir une adresse invalide
\end{itemize}

Un pointeur non initialisé ne pointe vers rien de fiable.

\subsection{Deux mondes : adresses vs valeurs}

Avec les pointeurs, il est essentiel de distinguer deux mondes :

\begin{itemize}
    \item le \textbf{monde des valeurs}
    \item le \textbf{monde des adresses}
\end{itemize}

La plupart des erreurs viennent de la confusion entre ces deux niveaux.

\begin{center}
\textbf{valeur pointée $\neq$ pointeur}
\end{center}

Le pointeur contient une adresse,
la valeur pointée se trouve ailleurs en mémoire.

Chaque fois que vous manipulez un pointeur,
vous devriez vous demander :

\begin{center}
\textit{Suis-je en train de manipuler une adresse ou une valeur ?}
\end{center}

\subsection{Allouer un pointeur (allocation dynamique)}

Dire que l’on ``alloue un pointeur'' est en réalité un abus de langage.

Ce que l’on fait réellement est :

\begin{enumerate}
    \item créer un objet quelque part en mémoire
    \item stocker son adresse dans le pointeur
\end{enumerate}

Dans l’analogie :

\begin{itemize}
    \item on construit une maison
    \item puis on écrit son adresse dans le carnet
\end{itemize}

L’allocation agit donc sur la \textbf{valeur pointée},
pas sur le pointeur lui-même.

\subsection{Libérer un pointeur}

Là encore, l’expression est trompeuse.

On ne libère pas le pointeur,
on libère la mémoire pointée.

\begin{itemize}
    \item la zone mémoire est rendue au système
    \item le pointeur conserve toujours l’adresse
\end{itemize}

Dans l’analogie :

\begin{itemize}
    \item on détruit la maison
    \item mais l’adresse reste écrite dans le carnet
\end{itemize}

Le pointeur devient alors \textbf{dangling}
(il contient une adresse invalide).

C’est pourquoi, après une libération,
il est recommandé de mettre le pointeur à \texttt{NULL}.

\subsection{Affectation de pointeurs}

Considérons l’instruction :

\begin{lstlisting}[language=C]
p1 = p2;
\end{lstlisting}

Cela signifie simplement :

\begin{center}
\textbf{copier une adresse}
\end{center}

On ne copie pas la valeur pointée,
on copie uniquement l’adresse.

Dans l’analogie :

\begin{itemize}
    \item quelqu’un vous donne une adresse
    \item vous la recopiez dans votre carnet
\end{itemize}

Les deux pointeurs désignent alors le même objet.

\subsection{Le pointeur nul : \texttt{NULL}}

L’instruction :

\begin{lstlisting}[language=C]
p = NULL;
\end{lstlisting}

signifie que le pointeur ne contient aucune adresse valide.

\begin{itemize}
    \item cela n’affecte pas la mémoire pointée auparavant
    \item cela efface seulement l’adresse stockée
\end{itemize}

Dans l’analogie :

\begin{itemize}
    \item on gomme la page du carnet
    \item la maison existe peut-être encore
\end{itemize}

Il est crucial de ne pas confondre :

\begin{itemize}
    \item \texttt{free(p)} : libère la mémoire pointée
    \item \texttt{p = NULL} : efface l’adresse
\end{itemize}

\section{C04c — Utiliser les pointeurs en pratique en C}

\subsection{Déclaration d’un pointeur}

Un pointeur est une variable, mais son type indique
la nature des données pointées.

La syntaxe générale est :

\begin{lstlisting}[language=C]
type *nom;
\end{lstlisting}

Exemple :

\begin{lstlisting}[language=C]
int *px;
\end{lstlisting}

Cela signifie :

\begin{itemize}
    \item \texttt{px} est un pointeur
    \item il pointe vers une zone mémoire contenant des \texttt{int}
\end{itemize}

Le type du pointeur indique donc comment lire la mémoire pointée.

\paragraph{Bonne pratique}

Comme toute variable, un pointeur devrait être initialisé.

\begin{lstlisting}[language=C]
int *px = NULL;
\end{lstlisting}

Cela garantit qu’il ne contient pas une adresse aléatoire.

\subsection{L’opérateur \& : obtenir une adresse}

L’opérateur \texttt{\&} permet de passer du monde des valeurs
au monde des adresses.

\begin{lstlisting}[language=C]
int x = 3;
int *px = &x;
\end{lstlisting}

Ici :

\begin{itemize}
    \item \texttt{\&x} retourne l’adresse de \texttt{x}
    \item cette adresse est stockée dans \texttt{px}
    \item \texttt{px} pointe maintenant vers \texttt{x}
\end{itemize}

On peut visualiser :

\begin{center}
\texttt{px} $\rightarrow$ \texttt{x} = 3
\end{center}

\subsection{L’opérateur * : accéder à la valeur pointée}

L’opérateur \texttt{*} permet de passer du monde des adresses
au monde des valeurs.

\begin{lstlisting}[language=C]
int x = 3;
int *px = &x;
printf("%d\n", *px);
\end{lstlisting}

\begin{itemize}
    \item \texttt{px} contient une adresse
    \item \texttt{*px} lit la valeur à cette adresse
\end{itemize}

L’instruction affiche donc :

\begin{center}
3
\end{center}

\paragraph{Conseil mental}

Au début, verbalisez :

\begin{center}
\textit{*px = contenu pointé par px (déréférencement)}
\end{center}
(lors de l'utilisation à la déclaration)\\
Cela aide à ne pas confondre adresse et valeur.

\subsection{Relation fondamentale}

Pour toute variable \texttt{x} :

\begin{lstlisting}[language=C]
*(&x) == x
\end{lstlisting}

Pourquoi ?

\begin{itemize}
    \item \texttt{\&x} donne l’adresse de \texttt{x}
    \item \texttt{*} lit le contenu à cette adresse
\end{itemize}

On retombe donc sur la valeur initiale.

\subsection{Modifier une valeur via un pointeur}

Un pointeur permet aussi de modifier directement la variable pointée.

\begin{lstlisting}[language=C]
int x = 3;
int *px = &x;

*px = 10;
\end{lstlisting}

Après cette instruction :

\begin{center}
\texttt{x == 10}
\end{center}

On a modifié la variable originale via son adresse.

\subsection{Pointeurs vers structures}

Si un pointeur pointe vers une structure :

\begin{lstlisting}[language=C]
struct Point {
    int x;
    int y;
};

struct Point p = {1,2};
struct Point *ptr = &p;
\end{lstlisting}

On peut accéder aux champs de deux façons :

\begin{lstlisting}[language=C]
(*ptr).x
\end{lstlisting}

ou, plus simplement :

\begin{lstlisting}[language=C]
ptr->x
\end{lstlisting}

La notation \texttt{->} est donc un raccourci pour :

\begin{center}
\texttt{(*ptr).champ}
\end{center}

\subsection{Attention au symbole *}

Le symbole \texttt{*} a plusieurs rôles en C :

\begin{itemize}
    \item multiplication
    \item déclaration de pointeur
    \item accès à la valeur pointée
\end{itemize}

Il faut bien distinguer :

\begin{itemize}
    \item \texttt{int *p;} \hspace{1em} déclaration
    \item \texttt{*p} \hspace{2.2em} accès mémoire
\end{itemize}

La déclaration n’est exécutée qu’une fois,
mais l’accès mémoire peut apparaître partout dans le programme.

\section{C04d — Pointeurs et fonctions : passage par ``référence''}

\subsection{Il n’existe pas de passage par référence en C}

En C, il n’existe que le passage par valeur.

Même lorsqu’on parle de ``passage par référence'',
on fait en réalité :

\begin{center}
\textbf{un passage par valeur de pointeur}
\end{center}

Pourquoi ?

\begin{itemize}
    \item un pointeur est une adresse
    \item si on passe une adresse à une fonction
    \item la fonction peut modifier le contenu à cette adresse
\end{itemize}

On obtient donc un effet équivalent au passage par référence.

\subsection{Exemple : \texttt{scanf}}

Considérons :

\begin{lstlisting}[language=C]
int x;
scanf("%d", &x);
\end{lstlisting}

Pourquoi met-on \texttt{\&x} ?

\begin{itemize}
    \item \texttt{scanf} doit modifier \texttt{x}
    \item pour modifier une variable, il faut connaître son adresse
    \item on passe donc l’adresse de \texttt{x}
\end{itemize}

Sans \texttt{\&}, la modification serait impossible.

\subsection{Exemple classique : \texttt{swap}}

Supposons :

\begin{lstlisting}[language=C]
int i = 3;
int j = 2;
swap(&i, &j);
\end{lstlisting}

On veut inverser les contenus mémoire de \texttt{i} et \texttt{j}.

La fonction correcte est :

\begin{lstlisting}[language=C]
void swap(int *a, int *b) {
    int tmp = *a;
    *a = *b;
    *b = tmp;
}
\end{lstlisting}

Les trois ingrédients du ``passage par référence'' :

\begin{enumerate}
    \item Les paramètres sont des pointeurs
    \item On utilise \texttt{*} pour accéder aux valeurs pointées
    \item On passe les arguments avec \texttt{\&}
\end{enumerate}

\subsection{Optimisation : éviter les copies}

Passer un objet volumineux par valeur peut être coûteux.

Exemple : matrices 100x100.

Solution :

\begin{itemize}
    \item passer un pointeur
    \item éviter la copie
\end{itemize}

Mais si la fonction ne doit pas modifier l’objet,
on protège avec \texttt{const} :

\begin{lstlisting}[language=C]
Complex add(const Complex *a, const Complex *b);
\end{lstlisting}

Cela signifie :

\begin{itemize}
    \item on passe des adresses (pas de copie)
    \item la fonction n’a pas le droit de modifier le contenu
\end{itemize}

\paragraph{Conseil pratique}

Quand vous passez un pointeur à une fonction :

\begin{center}
\textbf{mettez systématiquement \texttt{const} sauf si modification voulue}
\end{center}

\subsection{Retourner un pointeur : attention danger}

Considérons :

\begin{lstlisting}[language=C]
int* f() {
    int x = 5;
    return &x;
}
\end{lstlisting}

Ceci est \textbf{incorrect}.

Pourquoi ?

\begin{itemize}
    \item \texttt{x} est une variable locale
    \item sa durée de vie s’arrête à la fin de la fonction
    \item on retourne l’adresse d’une zone mémoire invalide
\end{itemize}

C’est un pointeur vers une variable qui n’existe plus.

\subsection{Solution : allocation dynamique}

Si on veut qu’un objet survive après la fin de la fonction,
il faut utiliser l’allocation dynamique :

\begin{lstlisting}[language=C]
int* f() {
    int *p = malloc(sizeof(int));
    *p = 5;
    return p;
}
\end{lstlisting}

Ici :

\begin{itemize}
    \item la mémoire est allouée dynamiquement
    \item sa durée de vie dépasse celle de la fonction
\end{itemize}

Mais :

\begin{center}
\textbf{tout malloc doit être accompagné d’un free}
\end{center}

\subsection{Erreur fréquente à éviter}

Ne jamais faire :

\begin{itemize}
    \item retourner l’adresse d’une variable locale
    \item oublier de libérer une mémoire allouée dynamiquement
\end{itemize}

Les pointeurs exigent une gestion rigoureuse
de la durée de vie des objets.

\section{C04e — Pointeurs et \texttt{const}}

\subsection{Règle générale pour \texttt{const}}

La règle pour comprendre \texttt{const} en C est simple :

\begin{center}
\textbf{\texttt{const} s’applique à ce qui précède,}
\textbf{sauf s’il n’y a rien avant, alors il s’applique à ce qui suit.}
\end{center}

Cette règle permet d’interpréter toutes les déclarations
impliquant des pointeurs et \texttt{const}.

\subsection{Pointeur vers donnée constante}

Considérons :

\begin{lstlisting}[language=C]
const int *p;
\end{lstlisting}

ou équivalent :

\begin{lstlisting}[language=C]
int const *p;
\end{lstlisting}

Ici, \texttt{const} s’applique au type \texttt{int}.

Cela signifie :

\begin{itemize}
    \item le pointeur peut changer d’adresse
    \item mais la valeur pointée ne peut pas être modifiée via ce pointeur
\end{itemize}

Autrement dit :

\begin{center}
\textbf{p est modifiable, mais *p est en lecture seule}
\end{center}

\subsection{Pointeur constant vers donnée modifiable}

Considérons :

\begin{lstlisting}[language=C]
int * const p;
\end{lstlisting}

Ici, \texttt{const} s’applique au pointeur lui-même.

Cela signifie :

\begin{itemize}
    \item le pointeur ne peut pas changer d’adresse
    \item mais la valeur pointée peut être modifiée
\end{itemize}

Autrement dit :

\begin{center}
\textbf{p est constant, mais *p est modifiable}
\end{center}

\subsection{Const sur les deux}

On peut combiner les deux :

\begin{lstlisting}[language=C]
const int * const p;
\end{lstlisting}

Cela signifie :

\begin{itemize}
    \item le pointeur est constant
    \item la valeur pointée est constante
\end{itemize}

On ne peut modifier ni l’adresse ni le contenu.

\subsection{Exemple détaillé}

Considérons :

\begin{lstlisting}[language=C]
int i = 2;
int j = 3;

const int *p1 = &i;
int * const p2 = &i;
\end{lstlisting}

\paragraph{Comportement de \texttt{p1}}

\begin{itemize}
    \item \texttt{*p1 = 5;} \hspace{1em} interdit
    \item \texttt{p1 = \&j;} \hspace{1.3em} autorisé
\end{itemize}

\texttt{p1} peut pointer ailleurs,
mais ne peut pas modifier le contenu pointé.

\paragraph{Comportement de \texttt{p2}}

\begin{itemize}
    \item \texttt{*p2 = 5;} \hspace{1em} autorisé
    \item \texttt{p2 = \&j;} \hspace{1.3em} interdit
\end{itemize}

\texttt{p2} pointe toujours vers la même adresse,
mais peut modifier la valeur pointée.

\subsection{Point subtil}

Le fait qu’un pointeur pointe vers un objet constant
ne signifie pas que cet objet ne peut jamais changer.

Cela signifie seulement :

\begin{center}
\textbf{cet objet ne peut pas être modifié au travers de ce pointeur}
\end{center}

Il peut être modifié par un autre pointeur non-const
ou directement via son nom.

\subsection{Interprétation mentale utile}

On peut résumer ainsi :

\begin{itemize}
    \item \texttt{const int *p} \hspace{1em} $\rightarrow$ pointeur vers lecture seule
    \item \texttt{int * const p} \hspace{1em} $\rightarrow$ pointeur figé
    \item \texttt{const int * const p} $\rightarrow$ tout figé
\end{itemize}

Comprendre cette distinction est essentiel
pour écrire des interfaces sûres
et éviter des modifications involontaires de données.

\section{C04f — Pointeurs vs références}

\subsection{En C : pas de références}

En C, il n’existe pas de références.

Le langage ne possède que :

\begin{itemize}
    \item des valeurs
    \item des pointeurs
\end{itemize}

Lorsque l’on parle de ``passage par référence'' en C,
il s’agit toujours d’un passage par valeur de pointeur.

Les références existent en revanche dans d’autres langages,
notamment en C++.

\subsection{Qu’est-ce qu’une référence ?}

Une référence peut être vue comme un \textbf{autre nom}
donné à une variable existante.

Contrairement à un pointeur :

\begin{itemize}
    \item une référence n’est pas une zone mémoire
    \item elle ne stocke pas une adresse
    \item elle est simplement un alias géré par le compilateur
\end{itemize}

On peut l’imaginer comme :

\begin{center}
\textbf{un drapeau posé sur une variable existante}
\end{center}

Le drapeau ne crée pas une nouvelle mémoire,
il désigne simplement une mémoire déjà existante.

\subsection{Différences fondamentales}

\paragraph{1. Existence en mémoire}

\begin{itemize}
    \item Un pointeur existe réellement en mémoire
    \item Une référence peut ne pas exister comme objet mémoire distinct
\end{itemize}

Un pointeur est une variable,
une référence est un alias.

\paragraph{2. Initialisation}

\begin{itemize}
    \item Un pointeur peut être non initialisé
    \item Une référence doit obligatoirement être initialisée
\end{itemize}

Une référence ne peut jamais être vide.

\paragraph{3. Valeur nulle}

\begin{itemize}
    \item Un pointeur peut valoir \texttt{NULL}
    \item Une référence ne peut jamais référencer ``rien''
\end{itemize}

Elle est toujours attachée à un objet valide.

\paragraph{4. Réaffectation}

\begin{itemize}
    \item Un pointeur peut être redirigé vers un autre objet
    \item Une référence ne peut jamais changer d’objet
\end{itemize}

Une référence désigne toujours le même objet.

\paragraph{5. Niveaux d’indirection}

\begin{itemize}
    \item On peut avoir des pointeurs vers pointeurs
    \item Les références ne s’enchaînent pas de cette manière
\end{itemize}

Les pointeurs permettent donc des structures beaucoup plus flexibles.

\subsection{Conséquences pratiques}

Les références sont plus simples à manipuler :

\begin{itemize}
    \item pas de gestion explicite d’adresses
    \item pas de pointeur nul
    \item moins d’erreurs possibles
\end{itemize}

C’est pourquoi, dans les langages qui offrent les deux :

\begin{center}
\textbf{on utilise des références quand c’est possible}
\end{center}

et

\begin{center}
\textbf{des pointeurs quand c’est nécessaire}
\end{center}

\subsection{Lien avec les usages des pointeurs}

On a vu que les pointeurs servent principalement à :

\begin{enumerate}
    \item partager un objet
    \item écrire du code générique
    \item gérer l’allocation dynamique
\end{enumerate}

Les références ne couvrent que le premier cas.

\begin{itemize}
    \item elles ne permettent pas la généricité via redirection
    \item elles ne permettent pas la gestion explicite de la mémoire
\end{itemize}

Les pointeurs restent donc indispensables en C.

\subsection{Idée finale}

On peut voir une référence comme :

\begin{center}
\textbf{un pointeur automatique que l’on ne peut pas modifier}
\end{center}

Tout est géré par le compilateur,
ce qui simplifie l’usage mais réduit la flexibilité.

En C, puisqu’il n’existe pas de références,
les pointeurs doivent être utilisés dans tous les cas
où l’on souhaite partager, manipuler ou contrôler des données en mémoire.
